\documentclass[11pt,twoside]{article}
\usepackage{etex}

\raggedbottom

%geometry (sets margin) and other useful packages
\usepackage{geometry}
\geometry{top=1in, left=1in,right=1in,bottom=1in}
 \usepackage{graphicx,booktabs,calc}
 
\usepackage{listings}


% Marginpar width
%Marginpar width
\newcommand{\pts}[1]{\marginpar{ \small\hspace{0pt} \textit{[#1]} } } 
\setlength{\marginparwidth}{.5in}
%\reversemarginpar
%\setlength{\marginparsep}{.02in}

 
%\usepackage{cmbright}lstinputlisting
%\usepackage[T1]{pbsi}


\usepackage{chngcntr,mathtools}
\counterwithin{figure}{section}
\numberwithin{equation}{section}

%\usepackage{listings}

%AMS-TeX packages
\usepackage{amssymb,amsmath,amsthm} 
\usepackage{bm}
\usepackage[mathscr]{eucal}
\usepackage{colortbl}
\usepackage{color}


\usepackage{subfigure,hyperref,enumerate,polynom,polynomial}
\usepackage{multirow,minitoc,fancybox,array,multicol}

\definecolor{slblue}{rgb}{0,.3,.62}
\hypersetup{
    colorlinks,%
    citecolor=blue,%
    filecolor=blue,%
    linkcolor=blue,
    urlcolor=slblue
}

%%%TIKZ
\usepackage{tikz}

\usepackage{pgfplots}
\pgfplotsset{compat=newest}

\usetikzlibrary{arrows,shapes,positioning}
\usetikzlibrary{decorations.markings}
\usetikzlibrary{shadows}
\usetikzlibrary{patterns}
%\usetikzlibrary{circuits.ee.IEC}
\usetikzlibrary{decorations.text}
% For Sagnac Picture
\usetikzlibrary{%
    decorations.pathreplacing,%
    decorations.pathmorphing%
}

\tikzstyle arrowstyle=[black,scale=2]
\tikzstyle directed=[postaction={decorate,decoration={markings,
    mark=at position .65 with {\arrow[arrowstyle]{stealth}}}}]
\tikzstyle reverse directed=[postaction={decorate,decoration={markings,
    mark=at position .65 with {\arrowreversed[arrowstyle]{stealth};}}}]
\tikzstyle dir=[postaction={decorate,decoration={markings,
    mark=at position .98 with {\arrow[arrowstyle]{latex}}}}]
\tikzstyle rev dir=[postaction={decorate,decoration={markings,
    mark=at position .98 with {\arrowreversed[arrowstyle]{latex};}}}]

\usepackage{ctable}

%
%Redefining sections as problems
%
\makeatletter
\newenvironment{exercise}{\@startsection 
	{section}
	{1}
	{-.2em}
	{-3.5ex plus -1ex minus -.2ex}
    	{1.3ex plus .2ex}
    	{\pagebreak[3]%forces pagebreak when space is small; use \eject for better results
	\large\bf\noindent{Exercise 1.\hspace{-1.5ex} }
	}
	}
	%{\vspace{1ex}\begin{center} \rule{0.3\linewidth}{.3pt}\end{center}}
	%\begin{center}\large\bf \ldots\ldots\ldots\end{center}}
\makeatother

%
%Fancy-header package to modify header/page numbering 
%
\usepackage{fancyhdr}
\pagestyle{fancy}
%\addtolength{\headwidth}{\marginparsep} %these change header-rule width
%\addtolength{\headwidth}{\marginparwidth}
%\fancyheadoffset{30pt}
%\fancyfootoffset{30pt}
\fancyhead[LO,RE]{\small Oke}
\fancyhead[RO,LE]{\small Page \thepage} 
\fancyfoot[RO,LE]{\small PS 2  } 
\fancyfoot[LO,RE]{\small \scshape CEE 260/MIE 273} 
\cfoot{} 
\renewcommand{\headrulewidth}{0.1pt} 
\renewcommand{\footrulewidth}{0.1pt}
%\setlength\voffset{-0.25in}
%\setlength\textheight{648pt}


\usepackage{paralist}

\newcommand{\osn}{\oldstylenums}
\newcommand{\lt}{\left}
\newcommand{\rt}{\right}
\newcommand{\pt}{\phantom}
\newcommand{\tf}{\therefore}
\newcommand{\?}{\stackrel{?}{=}}
\newcommand{\fr}{\frac}
\newcommand{\dfr}{\dfrac}
\newcommand{\ul}{\underline}
\newcommand{\tn}{\tabularnewline}
\newcommand{\nl}{\newline}
\newcommand\relph[1]{\mathrel{\phantom{#1}}}
\newcommand{\cm}{\checkmark}
\newcommand{\ol}{\overline}
\newcommand{\rd}{\color{red}}
\newcommand{\bl}{\color{blue}}
\newcommand{\pl}{\color{purple}}
\newcommand{\og}{\color{orange!90!black}}
\newcommand{\gr}{\color{green!40!black}}
\newcommand{\nin}{\noindent}
\newcommand{\la}{\lambda}
\renewcommand{\th}{\theta}
\newcommand*\circled[1]{\tikz[baseline=(char.base)]{
            \node[shape=circle,draw,thick,inner sep=1pt] (char) {\small #1};}}

\newcommand{\bc}{\begin{compactenum}[\quad--]}
\newcommand{\ec}{\end{compactenum}}

\newcommand{\n}{\\[2mm]}
%% GREEK LETTERS
\newcommand{\al}{\alpha}
\newcommand{\gam}{\gamma}
\newcommand{\eps}{\epsilon}
\newcommand{\sig}{\sigma}

\newcommand{\p}{\partial}
\newcommand{\pd}[2]{\frac{\partial{#1}}{\partial{#2}}}
\newcommand{\dpd}[2]{\dfrac{\partial{#1}}{\partial{#2}}}
\newcommand{\pdd}[2]{\frac{\partial^2{#1}}{\partial{#2}^2}}
\newcommand{\mr}{\mathbb{R}}
\newcommand{\xs}{x^{*}}
\newenvironment{solution}
{\medskip\par\quad\quad\begin{minipage}[c]{.8\textwidth}\gr}{\medskip\end{minipage}}

 
%%%%%%%%%%%%%%%%%%%%%%%%%%%%%%%%%%%%%%%%%%%%%%%%%%%
%%%%%%%%%%%%%%%%%%%%%%%%%%%%%%%%%%%%%%%%%%%%%%%%%%%

\begin{document}

\lstset{language=C++,
                basicstyle=\tiny\ttfamily,
                keywordstyle=\color{blue}\ttfamily,
                stringstyle=\color{red}\ttfamily,
                commentstyle=\color{gray}\ttfamily,
                morecomment=[l][\color{gray}]{\#}
}


\thispagestyle{empty}


\nin{\LARGE Problem Set 2 {\gr  }}\hfill{\bf Prof. Oke}

\medskip\hrule\medskip

\nin {\small CEE 260/MIE 273: Probability \& Statistics in Civil Engineering
\hfill\textit{ 09.11.2024}}

\bigskip

\nin{\it Due Tuesday, September 22, 2026 at 12:59 PM as PDF uploaded via Gradescope.
  If it helps and if possible, you can write your responses directly on this document and upload it instead.
  \textbf{Show as much work as possible in order to get FULL credit.}
  There FOUR problems %with a total of 30 points available.
}\\


\section*{Problem 1 \textit{(5 points)}}

Respond ``T'' ({\it True})  or  ``F'' (\textit{False}) to the following statements. Use the boxes provided. Each response is worth 1 point.

\begin{enumerate}[\bf (i)]
\item \hfill
  \begin{minipage}{.1\linewidth}
    \framebox(40,40){ \gr }
  \end{minipage}\quad
  \begin{minipage}{.85\linewidth}
    An impossible event has a probability less than or equal to 0. %False
  \end{minipage}
%  \underline{\hspace{10ex}} %F 

    \smallskip

% \item \hfill
%   \begin{minipage}{.1\linewidth}
%     \framebox(40,40){\gr T}
%   \end{minipage}\quad
%   \begin{minipage}{.85\linewidth}
%     Clustering is a supervised learning approach. %False
%   \end{minipage}

%   \smallskip
  
\item \hfill
  \begin{minipage}{.1\linewidth}
    \framebox(40,40){\gr  }
  \end{minipage}\quad
  \begin{minipage}{.85\linewidth}
    A certain event is equivalent to the sample space of outcomes.
  \end{minipage}

  \smallskip
  
\item \hfill
  \begin{minipage}{.1\linewidth}
    \framebox(40,40){\gr  }
  \end{minipage}\quad
  \begin{minipage}{.85\linewidth}
    If $A$ and $B$ are mutually exclusive, then $P(AB) \ne 0$.
  \end{minipage}

  \smallskip
  
\item \hfill
  \begin{minipage}{.1\linewidth}
    \framebox(40,40){\gr  }
  \end{minipage}\quad
  \begin{minipage}{.85\linewidth}
    Given an event $E$, then the probability of its complement can be given as $P(\ol{E}) = 1 + P(E)$
  \end{minipage}

  \smallskip

  \item \hfill
  \begin{minipage}{.1\linewidth}
    \framebox(40,40){\gr  }
  \end{minipage}\quad
  \begin{minipage}{.85\linewidth}
    $A$ and $B$ are collectively exhaustive events. If $P(B) = 0.4$, then $P(A) = 0.6$. %{\gr (Only true if $A$ and $B$ are mutually exclusive.)}
  \end{minipage}
  
%   \smallskip
  
% \item \hfill
%   \begin{minipage}{.1\linewidth}
%     \framebox(40,40){\gr  }
%   \end{minipage}\quad
%   \begin{minipage}{.85\linewidth}
%     The number of ways 6 books can be arranged on a bookshelf is 720. If two of the books are identical, then the total number of distinct arrangements is 360. %{\gr The permutations are $\fr{n!}{n_{1}!}$, where $n_{1} = 2$. Thus $\fr{6!}{2} = 360$.}
%   \end{minipage}
  
%   \smallskip
  
% \item \hfill
%   \begin{minipage}{.1\linewidth}
%     \framebox(40,40){\gr  }
%   \end{minipage}\quad
%   \begin{minipage}{.85\linewidth}
%     The number of distinct subgroups of size $n$ that can be formed from a larger group of $m$ objects is given by $\fr{m!}{n!(n-m)!}$.
%     %{\gr (The correct answer is $\fr{m!}{n!(m-n)!}$)}
%   \end{minipage}

%   \smallskip

%   \item \hfill
%   \begin{minipage}{.1\linewidth}
%     \framebox(40,40){\gr  }
%   \end{minipage}\quad
%   \begin{minipage}{.85\linewidth}
%     The events $E_{1}$ and $E_{2}$ are independent. If $P(E_{1}) = 0.4$ and $P(E_{1}E_{2}) = 0.04$, then $P(E_{2}) = 0.1$.
%     %{\gr ($P(E_{2}) = P(E_{1}E_{2})/P(E_{1}) = 0.04/0.4 = 0.1$)}
%   \end{minipage}
  \end{enumerate}

  % \smallskip
  
% \item \hfill
%   \begin{minipage}{.1\linewidth}
%     \framebox(40,40){}
%   \end{minipage}\quad
%   \begin{minipage}{.85\linewidth}
%     A good clustering solution should maximize the total within-cluster variation. %False
%   \end{minipage}

\eject  
\section*{Problem 2 \textit{(7 points)}}
Show your work/explanations in order to receive full credit.
Single number answers are not sufficient.

\begin{enumerate}[\bf (a)]

\item
% \item
%   \begin{minipage}{.8\linewidth}
    License plates in a certain state consist of 3 letters followed by 3 digits.\pts{2}
    How many different license plates can be manufactured?
    \vspace{16ex}
    
\item
  Your computer password must have only 6 characters.\pts{2} Each character can be a lowercase letter (e.g. \texttt{a}), an uppercase letter (e.g.
  \texttt{A}) or a digit (e.g. \texttt{0}). How many valid passwords could you possibly create?

  \vspace{17ex}
  
\item Your operating system is tightening its security protocols and asks you to update your password. \pts{3}
  It now requires that in addition to the earlier rules, you must now have at least one uppercase letter in your computer password.
  How many valid passwords could you possibly now create?


\vspace{20ex}
  
\end{enumerate}



% \item $P(A_3|E)P(E)$ \marginpar{\it[1pt]} % Bayes Theorem
  
%   \begin{tikzpicture}[scale=.8]
%     \draw[thick] (-4,-4) rectangle (7,3) node[above left] {$\bm S$};
%     \draw[thick] (1.5,0) ellipse  (4 cm and 2 cm) node[] {$\bm{E}$};
%     \draw (-1,3) -- (0,-4);
%     \draw (4,3) -- (4,-4);
%     \node at (-2,-3) {$\bm{A_1}$};
%     \node at (2,-3) {$\bm{A_2}$};
%     \node at (5,-3) {$\bm{A_3}$};
% %    \draw[thick] (3,0) circle (2 cm) node[right] {$\bm{E_2}$};
%  \end{tikzpicture}
 
 \section*{Problem 3 \textit{(8 points)}}
Data collected at elementary schools in DeKalb County, GA suggest that each
year roughly 25\% of students miss exactly one day of school, 15\% miss 2 days, and 30\% miss 3 or more days
due to sickness.

\begin{quote}
    \subsection*{\bl Example}
    \bl   Let $X$ be the number of school days missed in a year.
    
    The probability that a student chosen at random misses 2 or more days of school in a year is given by:
    \begin{equation*}
      P(X \ge 2) = P(X=2) + P(X \ge 3) = 0.15 + 0.30 = 0.45
  \end{equation*}
\end{quote}

\begin{enumerate}[\bf (a)]
\item  What is the probability that a student chosen at random does not miss any days of school due to sickness
  this year?
 \vspace{25ex}
\item What is the probability that a student chosen at random misses no more than one day?
 \vspace{25ex}

\item What is the probability that a student chosen at random misses at least one day?
 \vspace{25ex}

 \item What is the probability that a student chosen at random misses one or two days of school a year?
 \vspace{25ex}

 

\end{enumerate}

\section*{Problem 4 \textit{(MATLAB/Python}}

(Python version to be provided before midnight on Thursday)
\subsubsection*{ Objectives}
\begin{itemize}\item 
  Compute permutations and combinations

\item Learn how to import datasets

\item Compute probabilities

\item Create descriptive plots of datasets

\item Compute summary statistics 
\end{itemize}


\subsubsection*{Requirements}
You must have a working installation of MATLAB, preferably the most recent
version. Please edit and run your \texttt{.m} file in the same
directory as the supplied data (which you will download from Moodle). This
way, you will only need to use the filenames when loading data into MATLAB.

\subsubsection*{Submission instructions}
Use the template provided to write your code. Feel free to modify to your
liking. Append your name (\texttt{LASTNAME\_FIRSTNAME\_}) to \texttt{HW1.m} with an underscore (deleting
\texttt{\_template} from the filename) and then submit. Note that MATLAB
filenames must begin with a letter and contain only letters, numbers and
underscores. Any data or resources provided for an assignment will be available
to the graders so you only need to upload an \texttt{.m} file.

% \eject

% \begin{exercise}{Tutorial}
%   \textit{Optional}\quad If you have never used MATLAB before, visit this link:
%   \url{https://www.mathworks.com/help/matlab/learn_matlab/desktop.html}
%   and start with the lesson ``Desktop Basics.'' 
 
% \end{exercise}


\begin{exercise}{Summary statistics and probabilities}
The files \texttt{2013Jan\_Baltimore\_MeanTempF.csv} and
\texttt{2014Jan\_Baltimore\_MeanTempF.csv} (contained in the archive \texttt{HW1data.zip} contain average daily temperatures (in Fahrenheit)
recorded in Baltimore during the months January 2013 and January 2014,
respectively. 

\begin{enumerate}[(a)]
\item Use a built-in MATLAB function (\texttt{csvread} is a good one) to read
  the January 2013 values into an array $X$ and the January 2014 values into an
  array $Y$. (The first task has been completed in the template.)

\item Replace \texttt{??} in line 20 with the appropriate value to find the number of days in January 2013 with an average temperature of 30F. \pts{1}

\item Now uncomment the other lines to \pts{1} find the probability that the mean temperature in January 2013 was 30F and display your answer using the \texttt{fprintf} statement.

\item Write \pts{1} a one line statement in your script to compute the probability that the mean temperature in January 2013 was \textit{not} 30F. 

\item Display your answer to part (d) using \texttt{fprintf}. \pts{1}

\item Now compute the probability that the \pts{3} mean daily temperature in January 2014 was greater than 32F but \textit{no more} than 35F. Hint: to find the number of occurrences required, use:
  \begin{quote}
    \texttt{nnz(Y > 32 \& Y <= 35)}
  \end{quote}
Display your answer using \texttt{fprintf}.


  \item Create the following plots:
    \begin{enumerate}[\bf (i)]
    \item Histogram of \texttt{X} \pts{2}
    \item Histogram of \texttt{Y} \pts{2}
    \item Scatterplot of \texttt{Y} versus \texttt{X} \pts{2}.
  \end{enumerate}
  Make sure that in each plot, the $x$-axis and the $y$-axis are both labeled. Also be sure to include titles for each plot.
  When you or the TA runs your script, all three plots must be displayed simultaneously.
  Points will be deducted in each case where these instructions are not followed.
  Consult \texttt{example2.m} from Module 1 for examples of creating scatterplots and labeling axes and plots. \texttt{example1.m} shows how to plot histograms.

\item Use MATLAB \pts{4} to find the mean and standard deviation of the temperatures in both X and Y. Display your answers using two \texttt{fprintf} statements. (The first one has been commented out for you. Note the usage of square brackets when breaking up a long quote in MATLAB over multiple lines.)
  \end{enumerate}

% \begin{exercise}{Permutations and combinations}

%   \begin{enumerate}[\bf (a)]
%   \item Use the \texttt{perms} function to determine the number of ways you can rearrange \pts{2} the letters in the word ``MODEL''.
%     Be sure to uncomment the \texttt{fprintf} statement so that your answer is displayed when you run your script.

%   \item Now, find the number of \textit{distinct} ways you can rearrange \pts{2} the letters in the word ``LEVEL''.
%     Again, use an \texttt{fprintf} statement to display your answer.
%     (Hint: to find the number of unique rows in the output of \texttt{perms()}, use the following construction:
%     \begin{quote}
%       \texttt{size(unique(perms(`LEVEL'),"rows"))}
%     \end{quote}

%   \item Why is the answer \pts{2} in (b) different from that in (a)? Complete the \texttt{fprintf} statement given in the
%     template to explain your reasoning briefly. 

%   \item Given \pts{4} a Probability and Statistics class with 20 Civil/Environmental Engineering (CEE) majors, 10 Mechanical/Industrial Engineering (MIE) majors, and 5 students who are Other Engineering (OE) majors.
%     How many distinct groups  can you select from this class for a national engineering competition if the rules stipulate that all teams entering the competition must have exactly 10 members with 5 CEE students, 3 MIE students and 2 OE students?
%     Compute the answer to this problem using the \texttt{nchoosek(n,k)} function in MATLAB and \textit{display} your answer using \texttt{fprintf}.

%   \end{enumerate}
  
% \end{exercise}


% You are given the following probabilities of three events:
% \begin{align*}
%   P(A) &= 0.1 \\
%   P(B) &= 0.1 \\
%   P(B|A) &= 0.4 \\
%   P(C|AB) &= 0.2 \\
%   P(C|B) &= 0.4 
% \end{align*}
% Now, compute the following probabilities.

% \begin{enumerate}[\bf (a)]
% \item $P(AB)$ \pts{2}
%  \vspace{25ex}
 

% \item $P(A\cup B)$ \pts{2}
%  \vspace{25ex}
 

% \item $P(ABC)$ (Hint: this can be simplified as $P(C|AB)P(AB)$, using the multiplication rule) \pts{2}
%  \vspace{25ex}
 

% \eject
% \item $P(\ol{ABC})$ \pts{2}

%   % \begin{tikzpicture}[scale=.8]
% %     \draw[thick]  (-4,-3) rectangle (7,3); % node[above left] {$\bm S$};
% %     \draw[thick, pattern=north west lines,pattern color=green!50!black] (0,0) circle (2 cm) node[left] {$\bm{A}$};
% %     \draw[thick,  pattern=north west lines,pattern color=green!50!black] (3,0) circle (2 cm) node[right] {$\bm{B}$};
% %   \end{tikzpicture}

% %   % \begin{tikzpicture}[scale=.8]
% %   %   \draw[thick]  (-4,-3) rectangle (7,3); % node[above left] {$\bm S$};
% %   %   \draw[thick] (0,0) circle (2 cm) node[left] {$\bm{A}$};
% %   %   \draw[thick] (3,0) circle (2 cm) node[right] {$\bm{B}$};
% %   % \end{tikzpicture}

% %   \bigskip
  
% % \item $E_1\cap E_2$ \marginpar{\it[1]}
  
% %   \begin{tikzpicture}[scale=.8]
% %     \draw[thick]  (-4,-3) rectangle (7,3); % node[above left] {$\bm S$};
% %     \draw[thick] (0,0) circle (2 cm) node[left] {$\bm{E_1}$};
% %     \draw[thick] (3,0) circle (2 cm) node[right] {$\bm{E_2}$};

% %     \begin{scope}
% %       \clip (0,0) circle (2 cm) node[left] {$\bm{E_1}$};
% %       \fill[ pattern=north west lines,pattern color=green!50!black](3,0) circle (2 cm) node[right] {$\bm{E_2}$};
% %     \end{scope}
% %   \end{tikzpicture}


   
 
%   % \begin{tikzpicture}[scale=.8]
%   %   \draw[thick]  (-4,-5.5) rectangle (7,3); % node[above left] {$\bm S$};
%   %   \draw[thick] (0,0) circle (2 cm) node[left] {$\bm{A}$};
%   %   \draw[thick] (3,0) circle (2 cm) node[right] {$\bm{B}$};
%   %   \draw[thick] (1.5,-2.6) circle (2 cm) node[] {$\bm{C}$};

%   % \end{tikzpicture}
  
% \end{enumerate}
    
    
\end{document}

%%% Local Variables:
%%% mode: latex
%%% TeX-master: t
%%% End:
